\section{The classification problem and main result}
\label{sec:intro}

Consider the polynomial automorphism
\begin{equation}\label{eq:map}
 T(x,y,z)=(y,z,yz-x),\qquad
 T^{-1}(x,y,z)=(xy-z,x,y).
\end{equation}
Its integral dynamics preserve every level of
\begin{equation}\label{eq:invariant}
 K(x,y,z)=x^2+y^2+z^2-xyz,\qquad
 S_k(\Z)=\{P\in\Z^3:K(P)=k\},\quad k\in\Z.
\end{equation}
The time variable throughout is the ordinary iterate of the single map
$T$. In particular, we do not replace $T$ by an iterate, identify points
under coordinate sign changes, or ask for a finite orbit under every
trace-map group element.

Periodic curves of this map over the real or complex numbers are well
known. Their existence does not by itself determine the entire periodic
set on the integer lattice: a classification must exclude all other
periods, coordinate sizes, and levels at once. We give an elementary
arithmetic argument that does so. The crucial use of integrality occurs
after bounding the two neighbours of a coordinate of maximal modulus.
There are only finitely many remaining local alternatives, and their
resolution does not depend on the size of that maximum.

For $m\geq1$, put
\begin{equation}\label{eq:families}
 A_m=\Orb_T(m,0,0),\qquad
 B_m=\Orb_T(-1,m,-1),\qquad
 C_m=\Orb_T(1,m,1).
\end{equation}
Here $\Orb_T(P)=\{T^jP:j\in\Z\}$. We also use
\begin{equation}\label{eq:special}
 O=\{(0,0,0)\},\quad E=\{(2,2,2)\},\quad
 D=\{(-2,-2,2),(-2,2,-2),(2,-2,-2)\}.
\end{equation}
For a finite orbit, its height means the largest absolute scalar
coordinate at any point of that orbit.

\begin{theorem}[Complete integral periodic locus]\label{thm:main}
The periodic set of \eqref{eq:map} on $\Z^3$ is the disjoint union
\begin{equation}\label{eq:classification}
 \Per(T,\Z^3)=O\sqcup E\sqcup D\sqcup
 \bigsqcup_{m\geq1}(A_m\sqcup B_m\sqcup C_m).
\end{equation}
The orbit data are as in Table~\ref{tab:families}. Consequently the set
of least periods is exactly $\{1,3,4,6,12\}$.

An integral point has a bounded forward orbit, or a bounded backward
orbit, if and only if it belongs to \eqref{eq:classification}. Every
other integral point $P$ satisfies
\begin{equation}\label{eq:escape}
 \|T^nP\|_\infty\longrightarrow\infty
 \quad(n\to+\infty),\qquad
 \|T^nP\|_\infty\longrightarrow\infty
 \quad(n\to-\infty).
\end{equation}
\end{theorem}

\begin{table}[htbp]
\centering
\begin{tabular}{@{}lllll@{}}
\toprule
Orbit & Representative & Least period & Height & Level $K$\\
\midrule
$O$ & $(0,0,0)$ & $1$ & $0$ & $0$\\
$E$ & $(2,2,2)$ & $1$ & $2$ & $4$\\
$D$ & $(-2,-2,2)$ & $3$ & $2$ & $4$\\
$A_m$ & $(m,0,0)$ & $6$ & $m$ & $m^2$\\
$B_m$ & $(-1,m,-1)$ & $4$ & $m$ & $m^2-m+2$\\
$C_m$ & $(1,m,1)$ & $12$ & $m$ & $m^2-m+2$\\
\bottomrule
\end{tabular}
\caption{Complete integral orbit types. The parameter in the last three
rows ranges over all positive integers. Each row at each allowed
parameter denotes one orbit, not just a collection of representatives.}
\label{tab:families}
\end{table}

The escape statement in the theorem is proper escape on a discrete
lattice. It asserts neither monotonicity of the separate coordinates nor
a quantitative growth rate. The classification also yields a sharp
description on every level, even when $S_k(\Z)$ itself is infinite.

\begin{corollary}[Level support]\label{cor:support}
The total number of integral periodic points on $S_k(\Z)$ is $1$ when
$k=0$, $26$ when $k=4$, $6$ at every other positive square level,
$16$ at every other level $m^2-m+2$ with $m\geq1$, and $0$ otherwise.
\end{corollary}

We prove the classification in Sections~\ref{sec:itineraries}
and~\ref{sec:exhaustiveness}; the boundedness statement,
Corollary~\ref{cor:support}, and exact all-time return laws follow in
Section~\ref{sec:levels}.

\subsection{Classical inputs and the scope of the contribution}

In the half-trace coordinates $(X,Y,Z)=(x,y,z)/2$, the map is
$(X,Y,Z)\mapsto(Y,Z,2YZ-X)$ and the classical invariant is
\begin{equation}\label{eq:normalization}
 I=X^2+Y^2+Z^2-2XYZ-1=K/4-1.
\end{equation}
Thus our integral lattice becomes the half-integral lattice in that
normalization; a change of scale must be tracked in comparing arithmetic
statements. Roberts and Baake give the axis cycle, low-period families,
and their sign relations in their study of reversible trace maps
\citep[Eq.~(29), Table~I, and pp.~849--852]{roberts1994trace}.
In particular, the existence of the families in \eqref{eq:families}
and the relation between the four- and twelve-periodic families are not
new assertions. We verify their itineraries here because the exact
integral classification uses their least periods and signed members.

Roberts develops escape regions and growth criteria for trace maps over
the real and complex numbers \citep[\S\S4--5]{roberts1996escaping}.
Our argument does not improve those quantitative results. Instead, it
resolves the arithmetic alternatives at the complement and boundary of
possible escape: the equality case at modulus two, the one-zero branch,
the signed nonzero branch, and the unit cube. Its principal conclusion is
the exhaustiveness of \eqref{eq:classification}, not a new construction
of its displayed orbits.

Two other nearby questions have different quantifiers. Humphries
classifies points with finite orbit for the entire trace-map group
arising from $\Aut(F_2)$ \citep[Theorem~1]{humphries2016finite}.
Periodicity under one particular group element is weaker; the distinction
is illustrated in Section~\ref{sec:scope}. Ghosh and Sarnak study the
arithmetic of integral points on the surfaces \eqref{eq:invariant}
and their orbits under the Markoff group \citep{ghosh2022integral}.
A finite number of group orbits is not the assertion that each such
orbit is finite, and it does not enumerate cycles of the fixed word $T$.
Finally, the residual periodicity considered by Vishkautsan concerns a
two-reflection map on the zero Markoff surface and reduction modulo
primes \citep[\S\S1.1--1.2]{vishkautsan2016residual}; it is not the
all-level, single-map integral problem here. In particular, we do not
present the zero-level specialization alone as a new contribution.

These comparisons delimit the theorem rather than claim complete
literature priority. All arguments establishing
Theorem~\ref{thm:main} are included below and use only the recurrence,
integer arithmetic, and finiteness of bounded lattice subsets.
